% Standalone candidate dossier for the small-gap fourth-moment bootstrap.
% Candidate node, ledger acceptance pending: lem:mm-smallgap-fourth-moment is an
% orchestrator-accepted candidate id whose ledger entry has not yet been created.
\documentclass[../main.tex]{subfiles}
\begin{document}

% === SOLUTION HEADER =========================================================
% Certification is not recorded here: research/program/ledger.yaml owns the
% proofs[] mode and review path, and the review's front matter owns identities.
%
% ledger-node : lem:mm-smallgap-fourth-moment
% refines : lem:mm-smallgap-fourth-moment
% bounded_by : none (no bounded_by edge proposed for the candidate node; all
% author : claude-prover-w4p02
% date : 2026-08-30
% ============================================================================
\section*{Solution candidate: small-gap fourth-moment bootstrap from the published
frontier}
\label{sec:sol-lem-mm-smallgap-fourth-moment}

\paragraph{Refined statement and standing.}
The lemma is an \emph{unconditional implication}: it assumes a uniform Poincar\'e
constant $K_n$ for isotropic log-concave laws in dimension $n$ and bootstraps a
dimension-free fourth moment for first eigenfunctions whose gap is small on the scale
$1/K_n$.  Its intended instantiation is the \emph{published} frontier
$K_n=C_K\log n$ ($n\ge2$): Klartag \cite{Klartag2023Logarithmic} proves the Cheeger
constant bound $\psi_n\le C\sqrt{\log n}$, and Cheeger's inequality converts it into
$\CP\le C_K\log n$ for every isotropic log-concave probability on $\R^n$.  This
instantiation is the import node \texttt{thm:klartag-logn} (published import; ledger
acceptance pending as of this dossier's date); it is cited, not proved, here.  No
unreviewed preprint input occurs anywhere in this dossier; in particular
Theorem~\ref{thm:letwin-qcts} is not used.

\begin{hypothesis}[Uniform frontier input in dimension $n$]
\label{hyp:sol-sfm-frontier}
$K_n>0$ is a constant such that every isotropic log-concave probability measure on
$\R^n$ satisfies the Poincar\'e inequality with constant $K_n$: for every locally
Lipschitz $h\in L^2$,
\[
  \Var(h)\le K_n\,\E\abs{\nabla h}^2 .
\]
By the published import \texttt{thm:klartag-logn}
\cite{Klartag2023Logarithmic}, Hypothesis~\ref{hyp:sol-sfm-frontier} holds with
$K_n=C_K\log n$ for $n\ge2$, with $C_K$ universal.
\end{hypothesis}

\begin{theorem}[Small-gap fourth moment]
\label{thm:sol-lem-mm-smallgap-fourth-moment}
Assume Hypothesis~\ref{hyp:sol-sfm-frontier}.  Let
\[
  \dd\mu(x)=Z^{-1}e^{-V(x)}\dd x,
  \qquad V\in C^\infty(\R^n),\qquad\nabla^2V\succeq\varepsilon I_n,\ \varepsilon>0,
\]
be a smooth, strongly log-concave, centered \emph{isotropic} probability measure (a
regular isotropic approximant in the sense of Section~\ref{subsec:spectral-sde}), let
$L=\Delta-\nabla V\cdot\nabla$ be its Friedrichs realization, and let $f$ be a
normalized first nonconstant eigenfunction:
\[
  -Lf=\lambda f,\qquad\E_\mu f=0,\qquad\E_\mu f^2=1 .
\]
If
\[
  \lambda\ \le\ \frac{3}{8K_n},
\]
then
\begin{equation}
  \E_\mu f^4\ \le\ 2 .
  \label{eq:sol-sfm-main}
\end{equation}
The bound $2$ is universal: it depends on neither $n$, $\varepsilon$, $\mu$, nor
$\lambda$.  The regularization enters only through the qualitative finiteness
$\E_\mu f^4<\infty$.
\end{theorem}

\begin{proof}
\noindent\emph{Step 1: qualitative $L^4$ finiteness ($\varepsilon$-dependent, used only
for finiteness).}
The Bakry--\'Emery criterion \cite{BakryEmery1985} for the $\varepsilon$-uniformly
convex potential gives a logarithmic Sobolev inequality for $\mu$ with constant
depending only on $\varepsilon$, and hence, by Gross's theorem, hypercontractivity of
the semigroup $P_t=e^{tL}$: there are $t_0=t_0(\varepsilon)>0$ and an exponent
$q(t_0)\ge4$ with $\norm{P_{t_0}g}_{L^4(\mu)}\le\norm g_{L^2(\mu)}$ (see
\cite[Ch.~5]{BakryGentilLedoux2014}).  Since $P_{t_0}f=e^{-\lambda t_0}f$,
\[
  \norm f_{L^4(\mu)}=e^{\lambda t_0}\norm{P_{t_0}f}_{L^4(\mu)}
  \le e^{\lambda t_0}\norm f_{L^2(\mu)}<\infty .
\]
Thus $\E_\mu f^4<\infty$.  The bound is $\varepsilon$-dependent and is used only to
license the rearrangement in Step~4; no quantitative use is made of it.  We also
record that $f\in C^\infty(\R^n)$ by elliptic regularity for the equation
$\Delta f-\nabla V\cdot\nabla f=-\lambda f$ with smooth coefficients, and that
$f$ lies in the form domain $W^{1,2}(\mu)$ with
$\E_\mu\abs{\nabla f}^2=\langle f,-Lf\rangle=\lambda$.

\medskip
\noindent\emph{Step 2: the eigenvalue--energy identity
$\lambda\,\E f^4=3\,\E f^2\abs{\nabla f}^2$.}
For $k\in\mathbb N$ let $\varphi_k:\R\to\R$ be the truncated cubic
\[
  \varphi_k(s)=
  \begin{cases}
    s^3, & \abs s\le k,\\
    k^3+3k^2(s-k), & s>k,\\
    -k^3+3k^2(s+k), & s<-k,
  \end{cases}
\]
which is $C^1$, odd, nondecreasing, globally Lipschitz with constant $3k^2$, and
satisfies $\varphi_k(0)=0$, $\varphi_k'(s)=3(\abs s\wedge k)^2$.  Since the Dirichlet
form of the Friedrichs realization is a Dirichlet form, composition of the
form-domain element $f$ with a Lipschitz function vanishing at $0$ stays in the form
domain, with the chain rule $\nabla\varphi_k(f)=\varphi_k'(f)\nabla f$; concretely,
$\varphi_k(f)\in L^2(\mu)$ because $\abs{\varphi_k(f)}\le3k^2\abs f$, and
$\abs{\nabla\varphi_k(f)}\le3k^2\abs{\nabla f}\in L^2(\mu)$.  Pairing the
eigenequation with $\varphi_k(f)$ in the form,
\begin{equation}
  \lambda\,\E_\mu\bigl[f\,\varphi_k(f)\bigr]
  =\E_\mu\bigl[\nabla f\cdot\nabla\varphi_k(f)\bigr]
  =\E_\mu\bigl[\varphi_k'(f)\abs{\nabla f}^2\bigr]
  =3\,\E_\mu\bigl[(\abs f\wedge k)^2\abs{\nabla f}^2\bigr].
  \label{eq:sol-sfm-truncated}
\end{equation}
Both extreme members are monotone nondecreasing in $k$:
$s\,\varphi_k(s)\ge0$ increases pointwise to $s^4$ (for $\abs s>k$,
$s\varphi_k(s)=\abs s\,(k^3+3k^2(\abs s-k))\le s^4$ and is nondecreasing in $k$), and
$(\abs f\wedge k)^2\abs{\nabla f}^2$ increases pointwise to $f^2\abs{\nabla f}^2$.
Monotone convergence on both sides of \eqref{eq:sol-sfm-truncated} yields the identity
in $[0,\infty]$:
\begin{equation}
  \lambda\,\E_\mu f^4=3\,\E_\mu\bigl[f^2\abs{\nabla f}^2\bigr].
  \label{eq:sol-sfm-identity}
\end{equation}
By Step~1 the left side is finite, hence so is
$\E_\mu[f^2\abs{\nabla f}^2]=\lambda\,\E_\mu f^4/3$ (recall $\lambda>0$: a first
nonconstant eigenfunction has $\lambda=\E\abs{\nabla f}^2>0$).

\medskip
\noindent\emph{Step 3: Poincar\'e applied to $f^2$.}
The measure $\mu$ is isotropic and log-concave, so
Hypothesis~\ref{hyp:sol-sfm-frontier} applies to it.  The function $h=f^2$ is $C^1$
(hence locally Lipschitz), lies in $L^2(\mu)$ by Step~1, and has
$\E\abs{\nabla h}^2=4\,\E[f^2\abs{\nabla f}^2]<\infty$ by Step~2.  Therefore
\begin{equation}
  \E_\mu f^4-1=\E_\mu f^4-\bigl(\E_\mu f^2\bigr)^2=\Var_\mu\bigl(f^2\bigr)
  \le K_n\,\E_\mu\abs{\nabla f^2}^2
  =4K_n\,\E_\mu\bigl[f^2\abs{\nabla f}^2\bigr].
  \label{eq:sol-sfm-poincare}
\end{equation}

\medskip
\noindent\emph{Step 4: closure of the bootstrap at $\lambda\le3/(8K_n)$.}
Substituting \eqref{eq:sol-sfm-identity} into \eqref{eq:sol-sfm-poincare} and using the
small-gap hypothesis,
\[
  \E_\mu f^4-1
  \le\frac{4K_n\lambda}{3}\,\E_\mu f^4
  \le\frac{4K_n}{3}\cdot\frac{3}{8K_n}\,\E_\mu f^4
  =\frac12\,\E_\mu f^4 .
\]
Since $\E_\mu f^4<\infty$ (Step~1), the term $\tfrac12\E_\mu f^4$ may be subtracted,
giving $\tfrac12\E_\mu f^4\le1$, i.e.\ \eqref{eq:sol-sfm-main}.
\end{proof}

\begin{remark}[The external frontier input is genuinely needed]
\label{rem:sol-sfm-no-self-bootstrap}
The trivially available Poincar\'e constant for $\mu$ itself is $K=\CP(\mu)=1/\lambda$
(sharp for the first eigenfunction).  With that choice the coefficient in Step~4 is
$4K\lambda/3=4/3>1$, and the bootstrap closes for no $L^p$, $p>2$: the analogous
computation (pairing the eigenequation with truncations of $\abs s^{p-2}s$, which gives
$\lambda\,\E\abs f^p=(p-1)\,\E[\abs f^{p-2}\abs{\nabla f}^2]$, then Poincar\'e for
$\abs f^{p/2}$) produces the coefficient $Kp^2\lambda/(4(p-1))$, which at $K=1/\lambda$
equals $p^2/(4(p-1))=1+(p-2)^2/(4(p-1))>1$ for every $p>2$, independently of the
dimension.  An \emph{external} bound $K_n$ with
$K_n\lambda$ small is therefore what makes the argument close, and only on the
small-gap branch $\lambda\le3/(8K_n)$.  This is the probe's structural observation,
recorded here as a remark; it is not a claim and no statement is derived from it.
This use of the frontier is not circular for the downstream assembly: the pipeline
consuming this lemma outputs $\CP=O(\log^2n)$, strictly \emph{weaker} than its input
$O(\log n)$; the import is a calibrated external ingredient, not a hidden assumption
of the conclusion.
\end{remark}

\begin{remark}[Large-gap branch]
\label{rem:sol-sfm-large-gap}
Measures whose first eigenvalue violates the small-gap hypothesis satisfy
$\CP=1/\lambda<8K_n/3$ outright, at the frontier scale, and need no fourth-moment
control.  The branch split is performed by the consuming assembly dossier, not here.
\end{remark}

\begin{remark}[Hypotheses actually used]
\label{rem:sol-sfm-hypotheses}
The proof uses: smoothness and $\varepsilon$-strong log-concavity of $\mu$ (only for
elliptic regularity and the qualitative hypercontractive $L^4$ bound of Step~1);
centering and isotropy of $\mu$ (only so that Hypothesis~\ref{hyp:sol-sfm-frontier}
applies to $\mu$ in Step~3); the eigenequation through the Friedrichs form pairing;
the Dirichlet-form contraction and chain-rule property for Lipschitz post-composition;
and Hypothesis~\ref{hyp:sol-sfm-frontier}, discharged by the published import
\texttt{thm:klartag-logn} \cite{Klartag2023Logarithmic}.  No stochastic localization,
no Letwin input, no numerical evidence, and no unreviewed premise is used.  Modulo
ledger acceptance of the published import node, the implication proved here is
unconditional.
\end{remark}

\paragraph{Obstructions respected.}
The candidate node carries no \texttt{bounded\_by} edge; the registered route fences
were checked individually.  The argument is static and cut-free: no slice, excess, or
localized isoperimetric estimate occurs (\texttt{obs:two-tail},
\texttt{obs:circularity}, \texttt{obs:rank-one-refuted}); no projection or radial test
is promoted to a tensor bound (\texttt{obs:proj-ceiling}); no covariance occupation
functional appears (\texttt{obs:crude-insufficient},
\texttt{obs:relative-ceiling}).  Regarding \texttt{obs:relative-ceiling} specifically:
the a~priori input here is the published frontier constant $K_n$, used to produce a
strictly weaker downstream output, not an all-measure relative occupation premise.  The
covariance-spike fence is untouched (no localization occurs).

% No \printbibliography here (matches modules/*.tex): main.tex prints the single shared
% bibliography. Standalone, \cite shows "[?]" and \ref shows "??" --- expected.
\end{document}
